\section{Introduction}
\label{sec:introduction}

%% Introduction, finalized 2026-09-13 after a full line-by-line editing pass
%% with the user. MOVES 1-3 are reproduced here too (previously pulled from
%% BReRLA_manuscript.tex 699-708) so the whole section lives in one file;
%% assemble.py's carry-over of those lines is now redundant but harmless.

The weapon-target assignment problem (WTA), first formalized by Manne in
1958~\cite{manne1958target}, remains an active research problem
today~\cite{ahuja2007exact} because of its continued relevance to modern
defense scenarios. WTA and its variants allocate weapons to targets to
minimize the targets' remaining values after engagement, and the problem is
known to be NP-hard even in its original static, single-stage
form~\cite{lloyd1986weapons}.

The dynamic variant studied here (DWTA) adds four features that make the
model more operationally realistic: (1) multiple engagement stages, (2)
weapon reload times, (3) ammunition limits, and (4) windows during which a
target can be engaged. Given the urgent and time-critical nature of
battlefield situations, the solution method itself must also be fast.

Existing methods are slow in different ways. Exact methods guarantee
optimality but are limited to small instances, making them impractical for
repeated use at every stage for realistically sized
problems~\cite{ahuja2007exact,kline2019weapon}. Heuristics are fast, but each
is typically designed around a specific set of constraints and does not
readily transfer to other settings~\cite{kolitz1988analysis,lee2003genetic}.
Learning-based approaches have been proposed to close this gap, but most
remain narrower in scope. Prior reinforcement learning methods target the
static problem, fix their network size to a particular instance size, or,
when they do scale, assign weapons one at a time, causing computational cost
to grow with the number of
weapons~\cite{na2023weapon,gaudet2023,wu2022dynamic,yoontransformer}.

Unlike single-stage (static) WTA, a DWTA engagement requires each weapon to
make two distinct decisions. First, it must decide whether to fire or reserve
its ammunition for a future engagement. This is a strategic participation
decision. Second, if it fires, it must decide which target to engage. Prior
learned approaches typically ask a single policy to make both decisions,
producing either a target assignment or a decision to skip firing for each
weapon within a stage~\cite{na2023weapon,gaudet2023,yoontransformer}. This is
consistent with the standard design in neural combinatorial
optimization~\cite{kool2018attention,kwon2020pomo}, but it requires the
policy to learn two heterogeneous decisions, which makes the learning 
more difficult. The two decisions are fundamentally different. Target
assignment can be determined from information already known at the current
stage, such as remaining target values, kill probabilities, and current
engagement legality. In contrast, whether to fire depends on the consequences
of preserving ammunition for future stages, which have not yet occurred.
Treating both decisions as a single learning problem obscures this
distinction, and we show that it comes at a measurable cost.

Our claim is that only one of these two decisions needs to be learned. Once
the set of firing weapons is known, their assignment to targets can be solved
directly: a simple greedy rule on marginal value, requiring no training or
learned parameters, provably achieves at least half of the best possible assignment
value (Proposition~\ref{prop:auction-half}). Deciding which weapons should
fire admits no equivalent shortcut. Holding fire provides no benefit within
the current stage and pays off only through future engagements, so any rule
based solely on the current stage can be arbitrarily poor. We show that such
a rule can destroy substantially less target value than an optimal policy
would (Proposition~\ref{prop:myopic-gap}). Thus, we train the model whether to fire or hold and, argue that this minimal learning can achieve reasonably good solutions.

We build this idea into a simple system. A graph neural network considers all
weapons and targets jointly and outputs one number per weapon: the
probability that it should fire in the current stage. All weapons are
decided in a single pass rather than one at a time. A greedy marginal-value rule then
assigns the firing weapons to targets, and a lightweight search revisits the
completed schedule by flipping individual fire-or-hold decisions and
retaining only changes that improve the result. Trained on small instances
(5 to 7 weapons, targets, and stages), the policy generalizes without
further training to problems up to ten times larger, reaching 70 weapons,
100 targets, and 15 stages.

The main contributions of this study are summarized as follows:

\begin{itemize}
\item \textbf{A decomposition with a reason for every part}: We split the
per-stage decision in two. Target assignment is decided by a greedy marginal-value rule with
a provable $\tfrac12$ approximation guarantee (Proposition~\ref{prop:auction-half}). It holds even when weapons decide independently and simultaneously
(Remark~\ref{rem:auction-kills-pileon}). Only the fire and hold (participation decision) is learned for each weapon within the stage (Proposition~\ref{prop:myopic-gap}).

\item \textbf{Learning strictly less, and outperforms baselines}: The
policy outputs one Bernoulli parameter per weapon per stage, not a target
distribution. This reduces the mean objective by 14.9\% across twelve
configurations and outperforms every richer alternative we trained.

\item \textbf{A provably safe improvement step}: At inference we flip individual fire-or-hold decisions and keep only those the re-simulated objective confirms. This reduces the mean objective from 0.1219 to 0.1179 at a budget of ten re-simulations and to 0.1155 at thirty, uses no learned parameters, and cannot make a schedule worse at any budget (Proposition~\ref{prop:search-monotone}).

\item \textbf{Simultaneous decisions with no coordination machinery}: All
weapons are decided in a single pass. The assignment guarantee holds for any
firing set, including one produced by weapons deciding independently
(Remark~\ref{rem:auction-kills-pileon}), so the contention that communication
layers and conflict handlers exist to resolve never arises. Adding
weapon-to-weapon attention under an identical recipe makes the pipeline
slightly worse and three times less stable across seeds
(Section~\ref{sec:ablation}).

\item \textbf{Rigorous, standardized benchmarking}: Evaluation spans four
size tiers up to $70\times100\times15$ under resource scarcity, against exact solution approach,
a greedy heuristic, the auction alone, a sequential decoder, and other RL approaches such as AM and POMO
baselines~\cite{kool2018attention,kwon2020pomo}, all under an identical
protocol.
\end{itemize}
